\documentclass[
]{jss}

%% recommended packages
\usepackage{orcidlink,thumbpdf,lmodern}

\usepackage[utf8]{inputenc}

\author{
Paul Govan~\orcidlink{0000-0002-1821-8492}\\GE Aerospace
}
\title{\pkg{PRA}: Quantitative Project Risk Analysis in \proglang{R}}

\Plainauthor{Paul Govan}
\Plaintitle{PRA: Quantitative Project Risk Analysis in R}
\Shorttitle{\pkg{PRA}: Project Risk Analysis}


\Abstract{
The \pkg{PRA} package for \proglang{R} provides a unified, open-source
implementation of the core analytical methods used in quantitative
project risk analysis: the Second Moment Method (SMM), Monte Carlo
simulation (MCS), earned value management (EVM), contingency and
sensitivity analysis, Bayesian risk inference and updating, sigmoidal
learning curves, and design structure matrices (DSM). It operationalizes
causal-network and structural-network frameworks for project risk in a
dependency-light \proglang{R} package available on CRAN. Version 0.5.0
further adds an experimental probabilistic network module for forward
simulation and Bayesian conditioning along causal chains, and exposes
every analytical function as a tool over the Model Context Protocol
(MCP), so that large-language-model agents can run project risk analyses
through a documented set of skills. The package adheres to rOpenSci
Statistical Software Review standards and is supported by a
continuous-integration test suite running across macOS, Windows, and
Ubuntu.
}

\Keywords{project risk analysis, Monte Carlo simulation, earned value
management, Bayesian inference, design structure matrix, \proglang{R}}
\Plainkeywords{project risk analysis, Monte Carlo simulation, earned
value management, Bayesian inference, design structure matrix, R}

%% publication information
%% \Volume{50}
%% \Issue{9}
%% \Month{June}
%% \Year{2012}
%% \Submitdate{}
%% \Acceptdate{2012-06-04}

\Address{
  }


% tightlist command for lists without linebreak
\providecommand{\tightlist}{%
  \setlength{\itemsep}{0pt}\setlength{\parskip}{0pt}}




\usepackage{amsmath} \usepackage{booktabs} \usepackage{longtable}

\begin{document}



\section{Introduction}\label{introduction}

Quantitative risk analysis is a foundational discipline in project
management, enabling organizations to assess schedule and cost
uncertainty, identify high-impact risk drivers, and set evidence-based
contingency reserves. This article presents \pkg{PRA}, an open-source
\proglang{R} package that unifies the core quantitative methods of
project risk analysis---schedule and cost uncertainty, earned value
management, sensitivity and contingency analysis, learning curves, and
dependency structure analysis---under a single, reproducible,
CRAN-validated API.

Despite broad industry adoption of quantitative methods, open-source
tooling for integrated project risk analysis in \proglang{R} remains
fragmented. Proprietary platforms such as Oracle Primavera Risk Analysis
and Palisade @RISK provide integrated workflows but impose licensing
costs and limit reproducibility. Within \proglang{R}, no single package
unifies Monte Carlo schedule risk, earned value management, Bayesian
risk inference, sigmoidal learning curves, and dependency structure
analysis under a common API. The \pkg{mc2d} package \citep{mc2d}
provides general-purpose two-dimensional Monte Carlo simulation but has
no project-specific constructs. The \pkg{sensitivity} package
\citep{sensitivity} implements variance-based indices for general
functions but does not integrate with schedule or cost models. EVM
metrics are absent from CRAN entirely. Bayesian packages such as
\pkg{rstan} \citep{rstan} and \pkg{brms} \citep{brms} operate at a high
level of generality and require substantial expertise to apply in
project risk contexts.

\pkg{PRA} addresses this gap. It operationalizes the causal-network and
structural-network frameworks introduced by Govan (2014) and Govan and
Damnjanovic (2016), who developed the first implementation of causal
networks for project risk analysis grounded in the resource-based
view---demonstrating that organizational capabilities and shared
resources, not just schedule uncertainty, are the structural pathways
through which risks propagate across project tasks---and extended by
Govan and Damnjanovic (2020) with structural network measures that
quantify the topology of project risk networks. Building on this
foundation, \pkg{PRA} provides a coherent, CRAN-validated framework of
ten analytical modules covering the full quantitative project risk
workflow---from uncertainty characterization and simulation to
performance monitoring and Bayesian updating. The package targets
project analysts who work in \proglang{R} and require reproducible,
auditable outputs compatible with standard reporting pipelines. It
adheres to rOpenSci Statistical Software Review (SRR) standards
\citep{rOpenSciStats2021} and uses minimal hard dependencies: \pkg{mc2d}
\citep{mc2d} for triangular distribution families and \pkg{minpack.lm}
\citep{Elzhov2016} for nonlinear least squares fitting. The optional
Model Context Protocol (MCP) server activates through
\code{Suggests}-only dependencies.

The following table summarizes how \pkg{PRA} compares with the most
closely related \proglang{R} packages and the dominant commercial
platform.

\begin{table}

\caption{\label{tab:comparison-table}Comparison of \pkg{PRA} with related \proglang{R} packages and commercial tools. ``Exp.'' = experimental feature; ``---'' = not supported.}
\centering
\begin{tabular}[t]{lccccc}
\toprule
Feature & PRA & mc2d & sensitivity & rstan/brms & @RISK\\
\midrule
MCS schedule risk & Yes & Partial & --- & --- & Yes\\
EVM (11 metrics) & Yes & --- & --- & --- & Yes\\
Bayesian updating & Yes & --- & --- & Yes & ---\\
Sigmoidal learning curves & Yes & --- & --- & --- & ---\\
Design structure matrices & Yes & --- & --- & --- & ---\\
\addlinespace
Probabilistic networks & Exp. & --- & --- & Yes & ---\\
MCP tool server & Yes & --- & --- & --- & ---\\
Open-source (CRAN) & Yes & Yes & Yes & Yes & ---\\
SRR-compliant & Yes & --- & --- & --- & ---\\
\bottomrule
\end{tabular}
\end{table}

The package is installed from CRAN in the standard way:

\begin{CodeChunk}
\begin{CodeInput}
R> install.packages("PRA")
R> library("PRA")
\end{CodeInput}
\end{CodeChunk}

Optional dependencies---\pkg{ellmer} \citep{ellmer}, \pkg{mcptools}, and
\pkg{jsonlite}---are only required for the MCP server,
\code{pra\_mcp\_server()}.

The remainder of this article describes each analytical module in turn,
with working code examples drawn from the same \proglang{R} session.
Section\textasciitilde{}\ref{sec-uncertainty} covers schedule and cost
uncertainty methods (SMM, MCS, contingency, sensitivity, and correlation
matrix construction). Section\textasciitilde{}\ref{sec-evm} describes
earned value management. Section\textasciitilde{}\ref{sec-learning}
describes sigmoidal learning curves.
Section\textasciitilde{}\ref{sec-dsm} covers design structure matrices.
Section\textasciitilde{}\ref{sec-network} introduces the experimental
probabilistic network module. Section\textasciitilde{}\ref{sec-agent}
describes the Model Context Protocol integration.
Section\textasciitilde{}\ref{sec-casestudy} presents an integrated case
study. Section\textasciitilde{}\ref{sec-summary} concludes.

\section{Schedule and Cost Uncertainty}\label{sec-uncertainty}

Two complementary methods are provided for quantifying schedule and cost
uncertainty: the Second Moment Method for rapid analytical estimates,
and Monte Carlo simulation for full distributional analysis. Contingency
extraction and sensitivity analysis complete the workflow.

We illustrate the workflow with \code{building\_project}, a bundled
example dataset describing a mid-rise commercial building with six work
packages, adapted from the construction-project examples in
\citet{Damnjanovic2020}.

\begin{CodeChunk}
\begin{CodeInput}
R> data("building_project", package = "PRA")
R> bp <- building_project
R> bp$task_names
\end{CodeInput}
\begin{CodeOutput}
[1] "Site Preparation & Earthworks"  "Foundations"                   
[3] "Structural Frame"               "Building Envelope"             
[5] "Mechanical/Electrical/Plumbing" "Interior Fit-out & Finishes"   
\end{CodeOutput}
\end{CodeChunk}

\subsection{Second Moment Method}\label{second-moment-method}

The Second Moment Method (SMM) propagates task-level means and variances
to project-level estimates using the Central Limit Theorem
\citep{Benjamin2000}. For a project with \(n\) tasks:
\[E[X] = \sum_{i=1}^{n} E[X_i], \quad
  \text{Var}(X) = \sum_{i=1}^{n}\text{Var}(X_i) + 2\sum_{i<j}\rho_{ij}\sigma_i\sigma_j\]
The six task durations are triangular, so we obtain their means and
variances in closed form and propagate them with \code{smm()}:

\begin{CodeChunk}
\begin{CodeInput}
R> tri_mean <- function(d) (d$a + d$b + d$c) / 3
R> tri_var  <- function(d) (d$a^2 + d$b^2 + d$c^2 -
+                          d$a * d$b - d$a * d$c - d$b * d$c) / 18
R> task_means <- vapply(bp$task_distributions, tri_mean, numeric(1))
R> task_vars  <- vapply(bp$task_distributions, tri_var,  numeric(1))
R> 
R> result <- smm(task_means, task_vars)   # independent tasks
R> cat("Total mean:", round(result$total_mean, 2), "weeks\n")
\end{CodeInput}
\begin{CodeOutput}
Total mean: 71.67 weeks
\end{CodeOutput}
\begin{CodeInput}
R> cat("Total SD:  ", round(result$total_std,  2), "weeks\n")
\end{CodeInput}
\begin{CodeOutput}
Total SD:   4.53 weeks
\end{CodeOutput}
\end{CodeChunk}

Positive correlation between tasks inflates the total variance;
\code{smm()} accepts a correlation matrix and adds the covariance terms
analytically:

\begin{CodeChunk}
\begin{CodeInput}
R> result_cor <- smm(task_means, task_vars, bp$cor_mat)
R> cat("Total SD (correlated):", round(result_cor$total_std, 2), "weeks\n")
\end{CodeInput}
\begin{CodeOutput}
Total SD (correlated): 6.59 weeks
\end{CodeOutput}
\end{CodeChunk}

SMM is useful when only first and second moments are available and a
rapid estimate is needed. When task distributions are non-normal or
accurate tail behavior is required, Monte Carlo simulation is preferred.

\subsection{Monte Carlo Simulation}\label{monte-carlo-simulation}

\code{mcs()} propagates uncertainty through a project network by drawing
samples from user-specified task distributions---triangular, normal, or
uniform---over a large number of iterations \citep{Vose2008}. By default
the tasks are sampled independently; an optional correlation matrix is
supported, though its current handling is approximate
(Section\textasciitilde{}\ref{sec-summary}).

\begin{CodeChunk}
\begin{CodeInput}
R> set.seed(42)
R> results <- mcs(10000, bp$task_distributions)
R> cat("Simulated mean:", round(results$total_mean, 2), "weeks\n")
\end{CodeInput}
\begin{CodeOutput}
Simulated mean: 71.72 weeks
\end{CodeOutput}
\begin{CodeInput}
R> cat("Simulated SD:  ", round(results$total_sd,   2), "weeks\n")
\end{CodeInput}
\begin{CodeOutput}
Simulated SD:   4.57 weeks
\end{CodeOutput}
\end{CodeChunk}

The \code{mcs()} function returns an S3 object of class \code{"mcs"}
with a \code{print()} method that summarizes percentiles. The full
simulated distribution is available in
\code{results\$total\_distribution}. On a modern laptop, 10,000
simulations of a six-task network complete in under one second;
increasing the count to 100,000 provides more precise tail estimates at
roughly ten-fold additional runtime. The default of 10,000 balances
accuracy and speed for most practical analyses.

\subsubsection{Numerical Validation}\label{numerical-validation}

For independent tasks, the Central Limit Theorem makes the SMM estimates
the exact moments of the total, so the Monte Carlo simulation should
reproduce them. The two agree to within Monte Carlo error, validating
the sampling engine:

\begin{CodeChunk}
\begin{CodeInput}
R> cat(sprintf("Mean:  SMM %.2f  vs  MCS %.2f\n",
+             result$total_mean, results$total_mean))
\end{CodeInput}
\begin{CodeOutput}
Mean:  SMM 71.67  vs  MCS 71.72
\end{CodeOutput}
\begin{CodeInput}
R> cat(sprintf("SD:    SMM %.2f  vs  MCS %.2f\n",
+             result$total_std,  results$total_sd))
\end{CodeInput}
\begin{CodeOutput}
SD:    SMM 4.53  vs  MCS 4.57
\end{CodeOutput}
\end{CodeChunk}

\begin{CodeChunk}
\begin{CodeInput}
R> hist(results$total_distribution,
+   breaks = 50, freq = FALSE,
+   main = "MCS: Total Project Duration",
+   xlab = "Duration (weeks)", col = "steelblue", border = "white"
+ )
R> lines(density(results$total_distribution), col = "tomato", lwd = 2)
R> abline(v = results$total_mean, lty = 2, lwd = 1.5)
\end{CodeInput}
\begin{figure}

{\centering \includegraphics[width=0.8\linewidth]{pra-jss_files/figure-latex/mcs-hist-1} 

}

\caption[Distribution of simulated total project durations]{Distribution of simulated total project durations. The dashed vertical line marks the mean.}\label{fig:mcs-hist}
\end{figure}
\end{CodeChunk}

\subsection{Correlation Matrix
Construction}\label{correlation-matrix-construction}

Both \code{smm()} and \code{mcs()} accept a correlation matrix to model
dependence between tasks. \code{cor\_matrix()} constructs a valid
positive-definite correlation matrix by empirically sampling from a set
of task distributions, providing a convenient way to derive the
correlation input from the same distributional assumptions used
elsewhere in the workflow.

\begin{CodeChunk}
\begin{CodeInput}
R> dists <- list(
+   normal  = function(n) rnorm(n, mean = 10, sd = 2),
+   triang  = function(n) mc2d::rpert(n, min = 5, mode = 10, max = 15),
+   uniform = function(n) runif(n, min = 8, max = 12)
+ )
R> C_emp <- cor_matrix(num_samples = 1000, num_vars = 3, dists = dists)
R> print(round(C_emp, 3))
\end{CodeInput}
\begin{CodeOutput}
       [,1]   [,2]   [,3]
[1,]  1.000 -0.028 -0.022
[2,] -0.028  1.000 -0.010
[3,] -0.022 -0.010  1.000
\end{CodeOutput}
\end{CodeChunk}

\subsection{Contingency Analysis}\label{contingency-analysis}

\code{contingency()} extracts a reserve buffer as the difference between
a high-confidence percentile and the base (median) estimate---a standard
approach in project risk practice \citep{PMBOK2021}.

\begin{CodeChunk}
\begin{CodeInput}
R> reserve <- contingency(results, phigh = 0.80, pbase = 0.50)
R> cat("Schedule contingency (P80 - P50):", round(reserve, 2), "weeks\n")
\end{CodeInput}
\begin{CodeOutput}
Schedule contingency (P80 - P50): 3.98 weeks
\end{CodeOutput}
\end{CodeChunk}

\subsection{Sensitivity Analysis}\label{sensitivity-analysis}

\code{sensitivity()} decomposes total project variance by task
contribution, producing the inputs for a Tornado chart. Each task's
index is its own variance plus its covariance with every other task,
expressed as a proportion of total project variance.

\begin{CodeChunk}
\begin{CodeInput}
R> sens <- sensitivity(bp$task_distributions)
R> names(sens) <- paste0("T", seq_along(sens))
R> barplot(sort(sens),
+   horiz = TRUE, col = "steelblue",
+   main = "Sensitivity: Task Contributions",
+   xlab = "Proportion of Total Variance"
+ )
\end{CodeInput}
\begin{figure}

{\centering \includegraphics[width=0.8\linewidth]{pra-jss_files/figure-latex/sensitivity-1} 

}

\caption[Tornado chart]{Tornado chart: each bar shows a task's share of total project duration variance.}\label{fig:sensitivity}
\end{figure}
\end{CodeChunk}

Tasks with larger bars are the dominant drivers of total risk;
mitigation resources should be directed there first.

\section{Earned Value Management}\label{sec-evm}

The EVM module implements the complete ANSI/EIA-748 suite of eleven
performance metrics \citep{Fleming2010}. The three core
quantities---Planned Value (PV), Earned Value (EV), and Actual Cost
(AC)---are computed first; all performance indices and forecasts follow
from them.

\begin{CodeChunk}
\begin{CodeInput}
R> bac         <- 500000                          # budget at completion ($)
R> schedule    <- c(0.10, 0.25, 0.50, 0.75, 1.00) # cumulative value curve
R> time_period <- 3
\end{CodeInput}
\end{CodeChunk}

\begin{CodeChunk}
\begin{CodeInput}
R> pv_val <- pv(bac, schedule, time_period)
R> ev_val <- ev(bac, actual_per_complete = 0.40)
R> ac_val <- ac(c(45000, 110000, 135000), time_period, cumulative = FALSE)
R> 
R> cat("PV: $", format(pv_val, big.mark = ","), "\n")
\end{CodeInput}
\begin{CodeOutput}
PV: $ 250,000 
\end{CodeOutput}
\begin{CodeInput}
R> cat("EV: $", format(ev_val, big.mark = ","), "\n")
\end{CodeInput}
\begin{CodeOutput}
EV: $ 2e+05 
\end{CodeOutput}
\begin{CodeInput}
R> cat("AC: $", format(ac_val, big.mark = ","), "\n")
\end{CodeInput}
\begin{CodeOutput}
AC: $ 290,000 
\end{CodeOutput}
\end{CodeChunk}

\begin{CodeChunk}
\begin{CodeInput}
R> spi_val  <- spi(ev_val, pv_val)
R> cpi_val  <- cpi(ev_val, ac_val)
R> sv_val   <- sv(ev_val,  pv_val)
R> cv_val   <- cv(ev_val,  ac_val)
R> 
R> cat("SPI:", round(spi_val, 3), "  CPI:", round(cpi_val, 3), "\n")
\end{CodeInput}
\begin{CodeOutput}
SPI: 0.8   CPI: 0.69 
\end{CodeOutput}
\begin{CodeInput}
R> cat("SV: $", format(round(sv_val), big.mark = ","),
+     "  CV: $", format(round(cv_val), big.mark = ","), "\n")
\end{CodeInput}
\begin{CodeOutput}
SV: $ -50,000   CV: $ -90,000 
\end{CodeOutput}
\end{CodeChunk}

Both SPI \(< 1\) and CPI \(< 1\) indicate the project is behind schedule
and over budget. The \code{eac()} function provides three completion
forecasts under different efficiency assumptions:

\begin{CodeChunk}
\begin{CodeInput}
R> cat("EAC (typical):  $", format(round(eac(bac, method = "typical",
+     cpi = cpi_val)), big.mark = ","), "\n")
\end{CodeInput}
\begin{CodeOutput}
EAC (typical):  $ 725,000 
\end{CodeOutput}
\begin{CodeInput}
R> cat("EAC (atypical): $", format(round(eac(bac, method = "atypical",
+     ac = ac_val, ev = ev_val)), big.mark = ","), "\n")
\end{CodeInput}
\begin{CodeOutput}
EAC (atypical): $ 590,000 
\end{CodeOutput}
\begin{CodeInput}
R> cat("EAC (combined): $", format(round(eac(bac, method = "combined",
+     cpi = cpi_val, ac = ac_val, ev = ev_val, spi = spi_val)),
+     big.mark = ","), "\n")
\end{CodeInput}
\begin{CodeOutput}
EAC (combined): $ 833,750 
\end{CodeOutput}
\end{CodeChunk}

The remaining EVM metrics are:

\begin{CodeChunk}
\begin{CodeInput}
R> eac_typ <- eac(bac, method = "typical", cpi = cpi_val)
R> cat("ETC:  $", format(round(etc(bac, ev_val, cpi = cpi_val)), big.mark = ","), "\n")
\end{CodeInput}
\begin{CodeOutput}
ETC:  $ 435,000 
\end{CodeOutput}
\begin{CodeInput}
R> cat("TCPI:", round(tcpi(bac, ev_val, ac_val), 3), "\n")
\end{CodeInput}
\begin{CodeOutput}
TCPI: 1.429 
\end{CodeOutput}
\begin{CodeInput}
R> cat("VAC:  $", format(round(vac(bac, eac_typ)), big.mark = ","), "\n")
\end{CodeInput}
\begin{CodeOutput}
VAC:  $ -225,000 
\end{CodeOutput}
\end{CodeChunk}

The To-Complete Performance Index (TCPI) reports the cost efficiency
required on all remaining work to finish within the original budget;
values substantially above 1.0 signal an unrealistic recovery target.

\section{Learning Curves}\label{sec-learning}

Sigmoidal (S-curve) models capture the characteristic pattern of project
progress: slow start, rapid acceleration, and plateau approaching 100\%
completion. \pkg{PRA} provides three model types---Logistic, Pearl, and
Gompertz---fitted via the Levenberg--Marquardt nonlinear least-squares
algorithm in \pkg{minpack.lm} \citep{Elzhov2016}.

\begin{CodeChunk}
\begin{CodeInput}
R> progress <- data.frame(
+   time       = 1:9,
+   completion = c(5, 15, 40, 60, 70, 75, 80, 85, 90)
+ )
\end{CodeInput}
\end{CodeChunk}

\begin{CodeChunk}
\begin{CodeInput}
R> fit_log  <- fit_sigmoidal(progress, "time", "completion", "logistic")
R> fit_gomp <- fit_sigmoidal(progress, "time", "completion", "gompertz")
R> 
R> cat("Logistic RSE: ", round(summary(fit_log)$sigma,  2), "\n")
\end{CodeInput}
\begin{CodeOutput}
Logistic RSE:  4.15 
\end{CodeOutput}
\begin{CodeInput}
R> cat("Gompertz RSE: ", round(summary(fit_gomp)$sigma, 2), "\n")
\end{CodeInput}
\begin{CodeOutput}
Gompertz RSE:  2.89 
\end{CodeOutput}
\end{CodeChunk}

\begin{CodeChunk}
\begin{CodeInput}
R> plot_sigmoidal(fit_log, progress, "time", "completion", "logistic",
+                conf_level = 0.95,
+                main = "Logistic Learning Curve",
+                xlab = "Week", ylab = "Completion (%)")
\end{CodeInput}
\begin{figure}

{\centering \includegraphics[width=0.8\linewidth]{pra-jss_files/figure-latex/sigmoidal-plot-1} 

}

\caption[Logistic learning curve with 95\% confidence band]{Logistic learning curve with 95\% confidence band. The shaded region widens as extrapolation extends beyond the observed data range.}\label{fig:sigmoidal-plot}
\end{figure}
\end{CodeChunk}

\code{predict\_sigmoidal()} generates forecasts with confidence bounds
at arbitrary future time points:

\begin{CodeChunk}
\begin{CodeInput}
R> preds <- predict_sigmoidal(fit_log, seq(10, 12, by = 0.5),
+                             "logistic", conf_level = 0.95)
R> print(round(preds, 1))
\end{CodeInput}
\begin{CodeOutput}
     x pred  lwr  upr
1 10.0 84.3 78.2 90.3
2 10.5 84.3 78.2 90.4
3 11.0 84.3 78.2 90.4
4 11.5 84.3 78.2 90.5
5 12.0 84.3 78.2 90.5
\end{CodeOutput}
\end{CodeChunk}

\section{Dependency Structure Analysis}\label{sec-dsm}

Building on Govan and Damnjanovic (2016, 2020) and the foundational DSM
literature \citep{Steward1981, Browning2001}, the DSM module quantifies
structural coupling between project tasks through shared resources and
risks.

\subsection{Design Structure Matrices}\label{design-structure-matrices}

The \textbf{Parent DSM} is derived from the Resource--Task matrix
\(\mathbf{S}\) (rows = resources, columns = tasks) as
\(\mathbf{P} = \mathbf{S}^\top\mathbf{S}\). Off-diagonal entry
\(P_{jk}\) counts the number of resources shared between tasks \(j\) and
\(k\)---the structural pathway through which disruptions propagate
\citep{Govan2016}.

\begin{CodeChunk}
\begin{CodeInput}
R> S <- matrix(c(
+   1, 0, 1, 0,
+   1, 1, 0, 0,
+   0, 1, 0, 1,
+   0, 0, 1, 1,
+   0, 1, 1, 0
+ ), nrow = 4, ncol = 5)
R> rownames(S) <- paste0("R", 1:4)
R> colnames(S) <- paste0("T", 1:5)
\end{CodeInput}
\end{CodeChunk}

\begin{CodeChunk}
\begin{CodeInput}
R> p <- parent_dsm(S)
R> plot(p)
\end{CodeInput}
\begin{figure}

{\centering \includegraphics[width=0.8\linewidth]{pra-jss_files/figure-latex/parent-dsm-1} 

}

\caption[Parent DSM heatmap]{Parent DSM heatmap. Darker cells indicate more shared resources and stronger coupling between task pairs.}\label{fig:parent-dsm}
\end{figure}
\end{CodeChunk}

The \textbf{Grandparent DSM} adds a risk layer via the Risk--Resource
matrix \(\mathbf{R}\), tracing dependencies from risks through resources
to tasks. Off-diagonal entry \(G_{jk}\) counts the number of risks
shared between task \(j\) and task \(k\) through the resource chain---a
direct operationalization of the structural network measures developed
in Govan and Damnjanovic (2020).

\begin{CodeChunk}
\begin{CodeInput}
R> R_mat <- matrix(c(
+   1, 0, 1,
+   1, 1, 0,
+   0, 1, 0,
+   0, 0, 1
+ ), nrow = 3, ncol = 4)
R> rownames(R_mat) <- paste0("Risk", 1:3)
R> colnames(R_mat) <- paste0("R", 1:4)
\end{CodeInput}
\end{CodeChunk}

\begin{CodeChunk}
\begin{CodeInput}
R> g <- grandparent_dsm(S, R_mat)
R> plot(g)
\end{CodeInput}
\begin{figure}

{\centering 
\includegraphics[width=0.8\linewidth]{pra-jss_files/figure-latex/grandparent-dsm-1} 

}

\caption[Grandparent DSM heatmap]{Grandparent DSM heatmap. Shows risk-based coupling: task pairs that share risk exposure through common resources tend to fail together.}\label{fig:grandparent-dsm}
\end{figure}
\end{CodeChunk}

Task pairs with high off-diagonal values in the Grandparent DSM
represent concentrated risk exposure; these pairs are prime candidates
for contingency buffering or resource decoupling.

\section{Probabilistic Networks}\label{sec-network}

\textbf{Experimental.} The probabilistic network module extends DSM
structural analysis by attaching probability distributions to each node,
enabling forward simulation of cost uncertainty and Bayesian
conditioning on observed evidence. A probabilistic network is a directed
acyclic graph in which nodes represent project variables (risks,
resources, tasks, and the project total) and edges encode conditional
dependencies \citep{Pearl2009, Govan2014}. Where a DSM counts shared
resources between task pairs, a probabilistic network propagates
distributions through the causal chain from risks to total project cost.

\subsection{Building a Network}\label{building-a-network}

\code{prob\_net()} constructs the network from a node data frame, a
links data frame, and a named list of distributions. Distributions may
be \code{normal}, \code{lognormal}, \code{uniform}, \code{discrete}, or
\code{conditional} (switching between two sub-distributions depending on
a discrete parent node); \code{aggregate} nodes sum their parents.

\begin{CodeChunk}
\begin{CodeInput}
R> nodes <- data.frame(
+   id = c("A","B","C","D","E","F","G","H","I"),
+   label = c("Risk-1","Risk-2","Resource-1","Resource-2","Resource-3",
+             "Task-1","Task-2","Task-3","Project"),
+   stringsAsFactors = FALSE
+ )
R> links <- data.frame(
+   source = c("A","B","C","D","E","F","G","H"),
+   target = c("C","D","F","G","H","I","I","I"),
+   stringsAsFactors = FALSE
+ )
R> distributions <- list(
+   A = list(type = "discrete", values = c(1, 0), probs = c(0.70, 0.30)),
+   B = list(type = "discrete", values = c(1, 0), probs = c(0.60, 0.40)),
+   C = list(type = "conditional", condition = "A",
+     true_dist  = list(type = "normal", mean = 30000, sd = 8000),
+     false_dist = list(type = "normal", mean = 15000, sd = 3000)),
+   D = list(type = "conditional", condition = "B",
+     true_dist  = list(type = "normal", mean = 80000, sd = 20000),
+     false_dist = list(type = "normal", mean = 50000, sd = 10000)),
+   E = list(type = "normal", mean = 20000, sd = 4000),
+   F = list(type = "aggregate", nodes = "C"),
+   G = list(type = "aggregate", nodes = "D"),
+   H = list(type = "aggregate", nodes = "E"),
+   I = list(type = "aggregate", nodes = c("F","G","H"))
+ )
R> net <- prob_net(nodes, links, distributions = distributions)
\end{CodeInput}
\end{CodeChunk}

\subsection{Forward Simulation}\label{forward-simulation}

\code{prob\_net\_sim()} draws samples from every node in topological
order and returns a data frame with one column per node.

\begin{CodeChunk}
\begin{CodeInput}
R> sim_net <- prob_net_sim(net, num_samples = 10000)
R> cat("Mean project cost: $", format(round(mean(sim_net$I)), big.mark = ","), "\n")
\end{CodeInput}
\begin{CodeOutput}
Mean project cost: $ 113,612 
\end{CodeOutput}
\begin{CodeInput}
R> cat("P80 project cost:  $", format(round(quantile(sim_net$I, 0.80)),
+     big.mark = ","), "\n")
\end{CodeInput}
\begin{CodeOutput}
P80 project cost:  $ 135,543 
\end{CodeOutput}
\end{CodeChunk}

\subsection{Learning: Conditioning on
Evidence}\label{learning-conditioning-on-evidence}

\code{prob\_net\_learn()} clamps observed nodes to fixed values and
re-simulates the downstream distribution---an observational update in
the sense of \citet{Pearl2009}. The following call fixes Risk-2 (node B)
to absent (value 0), ruling out Technical Complexity as a concern:

\begin{CodeChunk}
\begin{CodeInput}
R> learn_net <- prob_net_learn(net, observations = list(B = 0),
+                             num_samples = 10000)
R> 
R> prior_net_dens <- density(sim_net$I)
R> post_net_dens  <- density(learn_net$I)
R> 
R> plot(prior_net_dens, col = "steelblue", lwd = 2,
+      main = "Total Project Cost: Prior vs. Posterior",
+      xlab = "Cost ($)",
+      xlim = range(c(sim_net$I, learn_net$I)),
+      ylim = range(c(prior_net_dens$y, post_net_dens$y)))
R> lines(post_net_dens, col = "tomato", lwd = 2)
R> legend("topright",
+   legend = c("Prior (Risk-2 uncertain)", "Posterior (Risk-2 = 0)"),
+   col = c("steelblue", "tomato"), lwd = 2, bty = "n")
\end{CodeInput}
\begin{figure}

{\centering \includegraphics[width=0.8\linewidth]{pra-jss_files/figure-latex/net-learn-1} 

}

\caption[Total project cost before and after observing that Risk-2 (Technical Complexity) did not occur]{Total project cost before and after observing that Risk-2 (Technical Complexity) did not occur. Conditioning on Risk-2 = 0 eliminates the heavy upper tail and shifts the distribution leftward by roughly \$30,000 in expectation.}\label{fig:net-learn}
\end{figure}
\end{CodeChunk}

\subsection{Intervention: The
do-Operator}\label{intervention-the-do-operator}

Conditioning answers ``what if we \emph{observe} that a risk did not
occur.'' A distinct question is ``what if we \emph{act} to eliminate
it.'' \code{prob\_net\_update()} performs structural surgery on the
network---adding or removing edges and replacing node
distributions---which corresponds to the \emph{do}-operator in Pearl's
causal calculus \citep{Pearl2009}, as opposed to
\code{prob\_net\_learn()}, which conditions without altering the graph.
Suppose management funds a mitigation that decouples Resource-2 from
Risk-2: remove the Risk-2 \(\to\) Resource-2 edge and fix Resource-2 to
its baseline cost.

\begin{CodeChunk}
\begin{CodeInput}
R> inter_net <- prob_net_update(net,
+   remove_links = data.frame(source = "B", target = "D"),
+   update_distributions = list(
+     D = list(type = "normal", mean = 50000, sd = 10000)
+   )
+ )
R> do_net <- prob_net_sim(inter_net, num_samples = 10000)
R> 
R> cat("Prior mean:            $", format(round(mean(sim_net$I)),   big.mark = ","), "\n")
\end{CodeInput}
\begin{CodeOutput}
Prior mean:            $ 113,612 
\end{CodeOutput}
\begin{CodeInput}
R> cat("Seeing (Risk-2 = 0):   $", format(round(mean(learn_net$I)), big.mark = ","), "\n")
\end{CodeInput}
\begin{CodeOutput}
Seeing (Risk-2 = 0):   $ 95,729 
\end{CodeOutput}
\begin{CodeInput}
R> cat("Doing  (do(Risk-2=0)): $", format(round(mean(do_net$I)),    big.mark = ","), "\n")
\end{CodeInput}
\begin{CodeOutput}
Doing  (do(Risk-2=0)): $ 95,625 
\end{CodeOutput}
\end{CodeChunk}

Here \emph{seeing} (\code{prob\_net\_learn()}) and \emph{doing}
(\code{prob\_net\_update()}) coincide because Risk-2 is an independent
root cause with no shared upstream confounder, so conditioning on it and
intervening on it induce the same downstream distribution. The two
diverge when risks share a common cause---intervening severs the
upstream pathway while observing propagates back through it---which is
where the \emph{do}-operator becomes indispensable.

\subsection{Ranking Risks by
Intervention}\label{ranking-risks-by-intervention}

Repeating the intervention for each root risk and measuring the
reduction in project-cost variance yields a principled, network-level
mitigation ranking---a causal analog of the Tornado chart of
Section\textasciitilde{}\ref{sec-uncertainty}:

\begin{CodeChunk}
\begin{CodeInput}
R> base_var <- var(sim_net$I)
R> 
R> do_A <- prob_net_sim(prob_net_update(net,
+   remove_links = data.frame(source = "A", target = "C"),
+   update_distributions = list(C = list(type = "normal", mean = 15000, sd = 3000))),
+   num_samples = 10000)
R> 
R> importance <- c(
+   "Risk-1" = base_var - var(do_A$I),
+   "Risk-2" = base_var - var(do_net$I)
+ )
R> barplot(sort(importance), horiz = TRUE, col = "steelblue",
+   main = "Risk Importance: Variance Eliminated by do(risk = 0)",
+   xlab = "Reduction in Project-Cost Variance")
\end{CodeInput}
\begin{figure}

{\centering 
\includegraphics[width=0.8\linewidth]{pra-jss_files/figure-latex/net-importance-1} 

}

\caption[Risk importance for the probabilistic network]{Risk importance for the probabilistic network: the reduction in total project-cost variance achieved by intervening to eliminate each root risk. Taller bars indicate higher-priority mitigation targets.}\label{fig:net-importance}
\end{figure}
\end{CodeChunk}

The probabilistic network API is marked experimental and may evolve in
future releases.

\section{Model Context Protocol Integration}\label{sec-agent}

Version 0.5.0 of \pkg{PRA} exposes all analytical functions as tools
callable by large-language-model (LLM) agents through the Model Context
Protocol (MCP) \citep{MCP2024}. \code{pra\_mcp\_server()} starts a
server that advertises each tool---Monte Carlo simulation, EVM, Bayesian
risk, learning curves, DSM, and the rest---so that Claude, Claude Code,
and other MCP-compatible clients can invoke them as native tools. The
server is registered in Claude Code with a single shell command:

\begin{CodeChunk}
\begin{CodeInput}
R> claude mcp add -s project pra -- Rscript -e "PRA::pra_mcp_server()"
\end{CodeInput}
\end{CodeChunk}

For Claude Desktop, add the following block to the MCP configuration
file:

\begin{verbatim}
{
  "mcpServers": {
    "pra": { "command": "Rscript",
             "args": ["-e", "PRA::pra_mcp_server()"] }
  }
}
\end{verbatim}

The tool definitions are built once with \code{ellmer::tool()}
\citep{ellmer} and served by \pkg{mcptools}. Each tool validates its
JSON inputs and returns a plain-text summary for the LLM together with,
where applicable, a rich HTML result with inline plots for clients that
render it.

To help an agent decide \emph{when} and \emph{how} to call each tool,
\pkg{PRA} bundles a set of Agent Skills under
\code{system.file("skills", package = "PRA")}. Each skill is a
\code{SKILL.md} file with structured guidance---when to use the method,
which tool to call, the JSON argument format, and how to interpret the
result---covering Monte Carlo simulation, sensitivity and contingency,
EVM, Bayesian risk, learning curves, and dependency structure analysis,
plus a routing overview. Pairing the skills with the MCP server lets an
agent carry out a complete project risk workflow (for example,
\code{mcs\_tool} followed by \code{contingency\_tool} and
\code{sensitivity\_tool}) from a natural-language request, without
embedding an LLM in the package itself.

The MCP layer is strictly optional: every analytical function operates
without any optional dependency. Only \code{pra\_mcp\_server()} requires
\pkg{ellmer}, \pkg{mcptools}, and \pkg{jsonlite}.

\section{Illustrative Case Study}\label{sec-casestudy}

This section demonstrates an integrated, end-to-end project risk
analysis for the \code{building\_project} dataset introduced in
Section\textasciitilde{}\ref{sec-uncertainty}, progressing through
uncertainty quantification, risk-driver ranking, performance monitoring,
completion forecasting, Bayesian risk updating, dependency structure
analysis, and network cost propagation in a single \proglang{R} session.
The workflow reuses objects built in the preceding sections (for example
the closed-form task moments, the fitted learning curve, and the
probabilistic network), illustrating how the modules compose within one
analysis.

\subsection{Step 1: Characterize Task
Uncertainty}\label{step-1-characterize-task-uncertainty}

Combine the rapid analytical estimate from the Second Moment
Method---reusing the closed-form triangular moments \code{task\_means}
and \code{task\_vars} from Section\textasciitilde{}\ref{sec-uncertainty}
with the dataset correlation matrix---with the full simulated
distribution from Monte Carlo simulation:

\begin{CodeChunk}
\begin{CodeInput}
R> bp <- building_project
R> set.seed(42)
R> 
R> smm_cs <- smm(task_means, task_vars, bp$cor_mat)   # analytic, correlated
R> sim     <- mcs(10000, bp$task_distributions)        # full distribution
R> res     <- contingency(sim, phigh = 0.80, pbase = 0.50)
R> 
R> cat("SMM total (correlated): mean", round(smm_cs$total_mean, 1),
+     "SD", round(smm_cs$total_std, 2), "weeks\n")
\end{CodeInput}
\begin{CodeOutput}
SMM total (correlated): mean 71.7 SD 6.59 weeks
\end{CodeOutput}
\begin{CodeInput}
R> cat("P50 schedule:", round(quantile(sim$total_distribution, 0.50), 1), "weeks\n")
\end{CodeInput}
\begin{CodeOutput}
P50 schedule: 71.6 weeks
\end{CodeOutput}
\begin{CodeInput}
R> cat("P80 schedule:", round(quantile(sim$total_distribution, 0.80), 1), "weeks\n")
\end{CodeInput}
\begin{CodeOutput}
P80 schedule: 75.6 weeks
\end{CodeOutput}
\begin{CodeInput}
R> cat("Contingency (P80-P50):", round(res, 2), "weeks\n")
\end{CodeInput}
\begin{CodeOutput}
Contingency (P80-P50): 3.98 weeks
\end{CodeOutput}
\end{CodeChunk}

\subsection{Step 2: Rank Risk Drivers}\label{step-2-rank-risk-drivers}

\code{sensitivity()} ranks the tasks by their share of total schedule
variance, focusing mitigation on the dominant drivers:

\begin{CodeChunk}
\begin{CodeInput}
R> drivers <- sensitivity(bp$task_distributions)
R> names(drivers) <- bp$task_names
R> print(round(sort(drivers, decreasing = TRUE), 3))
\end{CodeInput}
\begin{CodeOutput}
              Structural Frame Mechanical/Electrical/Plumbing 
                         0.304                          0.206 
   Interior Fit-out & Finishes              Building Envelope 
                         0.206                          0.133 
 Site Preparation & Earthworks                    Foundations 
                         0.076                          0.076 
\end{CodeOutput}
\end{CodeChunk}

\subsection{Step 3: Monitor Earned Value at
Mid-Project}\label{step-3-monitor-earned-value-at-mid-project}

\begin{CodeChunk}
\begin{CodeInput}
R> pv_cs <- pv(bp$bac, bp$schedule, bp$time_period)
R> ev_cs <- ev(bp$bac, bp$actual_per_complete)
R> ac_cs <- ac(bp$actual_costs, bp$time_period, cumulative = FALSE)
R> 
R> cat("SPI:", round(spi(ev_cs, pv_cs), 3),
+     " CPI:", round(cpi(ev_cs, ac_cs), 3), "\n")
\end{CodeInput}
\begin{CodeOutput}
SPI: 0.889  CPI: 0.986 
\end{CodeOutput}
\begin{CodeInput}
R> cat("EAC (typical): $", format(round(eac(bp$bac,
+     method = "typical", cpi = cpi(ev_cs, ac_cs))), big.mark = ","), "\n")
\end{CodeInput}
\begin{CodeOutput}
EAC (typical): $ 8,625,000 
\end{CodeOutput}
\end{CodeChunk}

\subsection{Step 4: Forecast
Completion}\label{step-4-forecast-completion}

Reusing the logistic learning curve \code{fit\_log} fitted in
Section\textasciitilde{}\ref{sec-learning}, project percent-complete at
future weeks with confidence bounds:

\begin{CodeChunk}
\begin{CodeInput}
R> preds_cs <- predict_sigmoidal(fit_log, seq(10, 12, by = 1),
+                               "logistic", conf_level = 0.95)
R> print(round(preds_cs, 1))
\end{CodeInput}
\begin{CodeOutput}
   x pred  lwr  upr
1 10 84.3 78.2 90.3
2 11 84.3 78.2 90.4
3 12 84.3 78.2 90.5
\end{CodeOutput}
\end{CodeChunk}

\subsection{Step 5: Update Risk Probability After Site
Inspection}\label{step-5-update-risk-probability-after-site-inspection}

\begin{CodeChunk}
\begin{CodeInput}
R> prior_r <- risk_prob(bp$cause_probs, bp$risks_given_causes,
+                      bp$risks_given_not_causes)
R> post_r  <- risk_post_prob(bp$cause_probs, bp$risks_given_causes,
+                           bp$risks_given_not_causes, bp$observed_causes)
R> cat("Prior P(delay):", round(prior_r, 3),
+     "-> Posterior P(delay):", round(post_r, 3), "\n")
\end{CodeInput}
\begin{CodeOutput}
Prior P(delay): 0.578 -> Posterior P(delay): 0.797 
\end{CodeOutput}
\end{CodeChunk}

\subsection{Step 6: Identify Coupled Tasks via
DSM}\label{step-6-identify-coupled-tasks-via-dsm}

\begin{CodeChunk}
\begin{CodeInput}
R> g_cs <- grandparent_dsm(bp$resource_task, bp$risk_resource)
R> plot(g_cs)
\end{CodeInput}
\begin{figure}

{\centering \includegraphics[width=0.8\linewidth]{pra-jss_files/figure-latex/cs-dsm-1} 

}

\caption[Grandparent DSM for the case-study project, showing which task pairs share the most risk exposure through common resources]{Grandparent DSM for the case-study project, showing which task pairs share the most risk exposure through common resources.}\label{fig:cs-dsm}
\end{figure}
\end{CodeChunk}

The Grandparent DSM reveals which task pairs share risk exposure through
common resources, informing targeted contingency allocation and resource
decoupling decisions.

\subsection{Step 7: Propagate Cost Through the
Network}\label{step-7-propagate-cost-through-the-network}

Finally, use the probabilistic network \code{net} from
Section\textasciitilde{}\ref{sec-network} to propagate cost uncertainty
to the project total, then condition on the observation that Risk-2
(Technical Complexity) did not occur:

\begin{CodeChunk}
\begin{CodeInput}
R> cat("Prior mean project cost: $",
+     format(round(mean(sim_net$I)), big.mark = ","), "\n")
\end{CodeInput}
\begin{CodeOutput}
Prior mean project cost: $ 113,612 
\end{CodeOutput}
\begin{CodeInput}
R> cat("Prior P80 project cost:  $",
+     format(round(quantile(sim_net$I, 0.80)), big.mark = ","), "\n")
\end{CodeInput}
\begin{CodeOutput}
Prior P80 project cost:  $ 135,543 
\end{CodeOutput}
\begin{CodeInput}
R> cat("Posterior mean (Risk-2 = 0): $",
+     format(round(mean(learn_net$I)), big.mark = ","), "\n")
\end{CodeInput}
\begin{CodeOutput}
Posterior mean (Risk-2 = 0): $ 95,729 
\end{CodeOutput}
\end{CodeChunk}

Conditioning on the resolved risk lowers the expected project cost,
closing the loop from uncertainty characterization to evidence-based
updating.

\section{Summary}\label{sec-summary}

\pkg{PRA} is a CRAN-available \proglang{R} package that provides a
unified, open-source implementation of quantitative project risk
analysis methods. It operationalizes the causal network and structural
network frameworks developed in Govan (2014) and Govan and Damnjanovic
(2016, 2020) through ten integrated analytical modules: SMM, MCS,
contingency analysis, sensitivity analysis, EVM (11 metrics), Bayesian
risk inference and updating, cost risk distributions, sigmoidal learning
curves, correlation matrix construction, and design structure matrices.
An experimental eleventh module provides probabilistic networks with
forward simulation and Bayesian conditioning.

Version 0.5.0 exposes every analytical function as a tool over the Model
Context Protocol, complemented by a bundled set of Agent Skills, so that
LLM agents can drive the package without any embedded model. All
exported functions adhere to rOpenSci SRR standards, with a test suite
verifying correctness, edge-case handling, and numerical stability
across three operating systems. The package is licensed under MIT and
the source is available at \url{https://github.com/paulgovan/PRA};
documentation and a companion online book are at
\url{https://paulgovan.github.io/PRA/}.

\subsection{Limitations and Future
Work}\label{limitations-and-future-work}

Several limitations apply. First, \pkg{PRA} currently propagates
uncertainty at the task level and does not yet model schedule
logic---critical path, activity precedence, and resource leveling. Full
schedule-network risk analysis, in which schedule Monte Carlo runs over
an activity-precedence network (PERT/CPM) to identify the critical path
and criticality indices under uncertainty, is a planned future direction
that would extend \pkg{PRA} from task-level to network-level schedule
risk. Second, \code{mcs()} induces task correlation by applying a
Cholesky factor to the sampled values, which reproduces the target rank
correlation but distorts the marginal means; the examples in this
article therefore treat tasks as independent for Monte Carlo simulation
and use \code{smm()}, which handles correlation analytically, when
correlation matters. A copula-based implementation that preserves the
marginals is planned. Third, the probabilistic network module is
experimental: the API may change and no convergence diagnostics are
provided for the Monte Carlo estimators. Fourth, the package currently
accepts data as \proglang{R} vectors and lists; direct import from
project management file formats (Microsoft Project, Primavera P6, CSV
exports from Jira or Asana) is not supported and is planned for a future
release.

\section*{Computational details}\label{computational-details}
\addcontentsline{toc}{section}{Computational details}

The results in this article were produced using \proglang{R}
\citep{RCore} with the \pkg{PRA} package (version 0.5.0) and its hard
dependencies \pkg{mc2d} \citep{mc2d} and \pkg{minpack.lm}
\citep{Elzhov2016}. The optional Model Context Protocol server
additionally uses \pkg{ellmer} \citep{ellmer}, \pkg{mcptools}, and
\pkg{jsonlite}. \proglang{R} itself and all packages used are available
from the Comprehensive \proglang{R} Archive Network (CRAN) at
\url{https://CRAN.R-project.org/}. A random seed of 42 is set at the
start of the session so that the simulation output is reproducible.

\renewcommand\refname{References}
\bibliography{pra-jss.bib}



\end{document}
